\documentclass[12pt,parskip=half-, twoside=semi]{scrartcl}
\usepackage{mathtools}
\usepackage[T1]{fontenc}
\usepackage[utf8]{inputenc}
\usepackage{multicol}
\usepackage{cmbright} % police sans serif
\usepackage{enumitem}
\SetEnumitemKey{cols}
{
  before = \begin{multicols}{#1},
    after = \end{multicols}
    }
\setlist[enumerate]{label=\textbf{\alph*)}, itemsep=1em}
\usepackage[use-files]{xsim}
\xsimsetup{
    path=xsim,
    load-style=mytemplate,
    exercise/template=mytemplate,
}


\usepackage[french]{babel}
\usepackage{fontawesome5}
\usepackage{tabularray}
\usepackage{scrlayer-scrpage}
%\setlength{\headheight}{50pt}
%\lohead{Fiches d'exercices de drill}
%\rohead{Chapitre 3 - Diviseurs et multiples}
\pagestyle{empty}

%Début du document%%%%%%%%%%%%%%%%%%%%%%%%%%%%%%%%%%%%%%%%%
\begin{document}

\begin{center}
  \bfseries
  \begin{Huge}
    Chapitre 3 - Diviseurs et multiples
  \end{Huge}
  \smallbreak
  \begin{large}
    Niveau 1 \faIcon{star} - Fiche 1
  \end{large}
  \smallbreak
  \rule[10pt]{.5\linewidth}{1 pt}
\end{center}

\begin{exercise}
  Décompose les nombres suivants en un produit de facteurs premiers.
  \begin{enumerate}[cols=3]
      \item 18
      \item 24
      \item 12
      \item 30
      \item 45
      \item 36
  \end{enumerate}
\end{exercise}

\begin{exercise}
  Écris l'égalité euclidienne qui découle de cette  division.
  \begin{enumerate}[cols=3]
    \item $15 : 4$
    \item $25 : 6$
    \item $14 : 3$
    \item $32 : 5$
    \item $18 : 7$
    \item $27 : 8$
  \end{enumerate}
\end{exercise}

\begin{exercise}
  Calcule le PGCD des nombres suivants.
\begin{enumerate}[cols=3]
    \item 12 et 18
    \item 24 et 36
    \item 15 et 25
    \item 9 et 27
    \item 14 et 28
    \item 8 et 12
 \end{enumerate}
\end{exercise}

\begin{exercise}
  Calcule le PPCM des nombres suivants.
  \begin{enumerate}[cols=3]
    \item 4 et 6
    \item 8 et 12
    \item 5 et 10
    \item 7 et 14
    \item 9 et 15
    \item 10 et 20
  \end{enumerate}
\end{exercise}

\begin{exercise}
    Écris l'expression littérale générale qui correspond aux énoncés suivants (avec \(n\) pour représenter les naturels).
    \begin{enumerate}[cols=2]
        \item Deux multiples différents de 3.
        \item Un nombre impair.
        \item Trois fois un nombre pair.
        \item La moitié d'un multiple de 4. 
        \item Un multiple de 3 augmenté de 21.
        \item Un multiple de 3 diminé de 2.
    \end{enumerate}
  \end{exercise}

\newpage 


\begin{center}
  \bfseries
  \begin{Huge}
    Chapitre 3 - Diviseurs et multiples
  \end{Huge}
  \smallbreak
  \begin{large}
    Niveau 1 \faIcon{star} - Fiche 2
  \end{large}
  \smallbreak
  \rule[10pt]{.5\linewidth}{1 pt}
\end{center}
  

\setcounter{exercise}{0}
  \begin{exercise}
      Décompose les nombres suivants en un produit de facteurs premiers.

    \begin{enumerate}[cols=3]
        \item 15
        \item 28
        \item 20
        \item 60
        \item 63
        \item 42
    \end{enumerate}
  \end{exercise}
   
  \begin{exercise}
    Écris l'égalité euclidienne qui découle de cette
  division.
    \begin{enumerate}[cols=3]
      \item $19 : 9$
      \item $21 : 4$
      \item $31 : 6$
      \item $16 : 5$
      \item $33 : 7$
      \item $20 : 9$
    \end{enumerate}
  \end{exercise}
  
  \begin{exercise}
    Calcule le PGCD des nombres suivants.
  \begin{enumerate}[cols=3]
      \item \(20 \) et \( 30\)
      \item \(10 \) et \( 50\) 
      \item \(7 \) et \( 21\)
      \item \(16 \) et \( 24\)
      \item \(40 \) et \( 60\)
      \item \(3 \) et \( 9\)   
    \end{enumerate}
  \end{exercise}
  
  \begin{exercise}
    Calcule le PPCM des nombres suivants.
    \begin{enumerate}[cols=3]
      \item 12 et 18
      \item 14 et 28
      \item 6 et 9
      \item 16 et 24
      \item 20 et 30
      \item  15 et 25
    \end{enumerate}
  \end{exercise}
  
  \begin{exercise}
      Écris l'expression littérale générale qui correspond aux énoncés suivants (avec \(n\) pour représenter les naturels).
      \begin{enumerate}[cols=2]
          \item Deux nombres pairs consécutifs.
          \item Trois nombres impairs consécutifs.
          \item Un nombre impair soustrait de 1.
          \item La somme de deux multiples de 3.
          \item Un multiple de 11.
          \item Le double d'un nombre.
      \end{enumerate}
    \end{exercise}

\newpage
\begin{center}
  \bfseries
  \begin{Huge}
    Correctif
  \end{Huge}
  \smallbreak
  \begin{large}
    Niveau 1 \faIcon{star} - Fiche 1
  \end{large}
  \smallbreak
  \rule[10pt]{.5\linewidth}{1 pt}
\end{center}

\setcounter{exercise}{0}


      \begin{exercise}
          Décompose les nombres suivants en un produit de facteurs premiers.

        \begin{enumerate}[cols=3]
          \item 18 = $2 \cdot 3 \cdot 3$
          \item 24 = $2 \cdot 2 \cdot 2 \cdot 3$
          \item 12 = $2 \cdot 2 \cdot 3$
          \item 30 = $2 \cdot 3 \cdot 5$
          \item 45 = $3 \cdot 3 \cdot 5$
          \item 36 = $2 \cdot 2 \cdot 3 \cdot 3$
        \end{enumerate}
      \end{exercise}
      
      
      \begin{exercise}
        Écris l'égalité euclidienne qui découle de cette
      division.
          \begin{enumerate}[cols=3]
            \item $15 = 4 \cdot 3 + 3$
            \item $25 = 6 \cdot 4 + 1$
            \item $14 = 3 \cdot 4 + 2$
            \item $32 = 5 \cdot 6 + 2$
            \item $18 = 7 \cdot 2 + 4$
            \item $27 = 8 \cdot 3 + 3$
          \end{enumerate}
      \end{exercise}
      
      \begin{exercise}
        Calcule le PGCD des nombres suivants.
          \begin{enumerate}[cols=3]
              \item 6
              \item 12
              \item 5
              \item 9
              \item 14
              \item 4
          \end{enumerate}
      \end{exercise}
      
      \begin{exercise}
        Calcule le PPCM des nombres suivants.
          \begin{enumerate}[cols=3]
              \item 12
              \item 24
              \item 10
              \item 14
              \item 45
              \item 20
          \end{enumerate}
      \end{exercise}
      
      \begin{exercise}
          Écris l'expression littérale générale qui correspond aux énoncés suivants (avec \(n\) pour représenter les naturels).
          \begin{enumerate}[cols=2]
            \item \(3n\) et \(3n+6\)
            \item \(2n+1\)
            \item \(3 \cdot 2n = 6n\)
            \item \(4n : 2 = 2n \)
            \item \(3n+21\)
            \item \(3n-2\)
          \end{enumerate}
        \end{exercise}

\newpage 

\begin{center}
  \bfseries
  \begin{Huge}
    Correctif
  \end{Huge}
  \smallbreak
  \begin{large}
    Niveau 1 \faIcon{star} - Fiche 2
  \end{large}
  \smallbreak
  \rule[10pt]{.5\linewidth}{1 pt}
\end{center}

\setcounter{exercise}{0}
    
    \begin{exercise}
        Décompose les nombres suivants en un produit de facteurs premiers.

      \begin{enumerate}[cols=3]
        \item 15 = $3 \cdot 5$
        \item 28 = $2 \cdot 2 \cdot 7$
        \item 20 = $2 \cdot 2 \cdot 5$
        \item 60 = $2 \cdot 2 \cdot 3 \cdot 5$
        \item 63 = $3 \cdot 3 \cdot 7$
        \item 42 = $2 \cdot 3 \cdot 7$
      \end{enumerate}
    \end{exercise}
    
    
    \begin{exercise}
      Écris l'égalité euclidienne qui découle de cette
    division.
        \begin{enumerate}[cols=3]
          \item $19 = 9 \cdot 2 + 1$
          \item $21 = 4 \cdot 5 + 1$
          \item $31 = 6 \cdot 5 + 1$
          \item $16 = 5 \cdot 3 + 1$
          \item $33 = 7 \cdot 4 + 5$
          \item $20 = 9 \cdot 2 + 2$
        \end{enumerate}
    \end{exercise}
    
    \begin{exercise}
      Calcule le PGCD des nombres suivants.
        \begin{enumerate}[cols=3]
            \item 10
            \item 10
            \item 7
            \item 8
            \item 20
            \item 3
        \end{enumerate}
    \end{exercise}
    
    \begin{exercise}
      Calcule le PPCM des nombres suivants.
        \begin{enumerate}[cols=3]
            \item 36
            \item 28
            \item 18
            \item 48
            \item 30
            \item 75
        \end{enumerate}
    \end{exercise}
    
    \begin{exercise}
        Écris l'expression littérale générale qui correspond aux énoncés suivants (avec \(n\) pour représenter les naturels).
        \begin{enumerate}[cols=2]
            \item 2n+2 , 2n+4
            \item 2n+3, 2n+5, 2n+7
            \item 1 - 2n + 1
            \item 3n + 3n+6
            \item 11n
            \item 2n 
        \end{enumerate}
      \end{exercise}

\newpage

\begin{center}
  \bfseries
  \begin{Huge}
    Chapitre 3 - Diviseurs et multiples
  \end{Huge}
  \smallbreak
  \begin{large}
    Niveau 1 \faIcon{star} - Fiche 3
  \end{large}
  \smallbreak
  \rule[10pt]{.5\linewidth}{1 pt}
\end{center}

\setcounter{exercise}{0}


    \begin{exercise}
        Décompose les nombres suivants en un produit de facteurs premiers.

      \begin{enumerate}[cols=3]
          \item 54
          \item 56
          \item 72
          \item 90
          \item 48
          \item 70
      \end{enumerate}
    \end{exercise}
    
    \begin{exercise}
      Écris l'égalité euclidienne qui découle de cette
    division.
      \begin{enumerate}[cols=3]
        \item $30 : 8$
        \item $28 : 7$
        \item $35 : 6$
        \item $45 : 9$
        \item $24 : 5$
        \item $22 : 4$
      \end{enumerate}
    \end{exercise}
    
    \begin{exercise}
      Calcule le PGCD des nombres suivants.
    \begin{enumerate}[cols=3]
        \item 6 et 9
        \item 25 et 35
        \item 11 et 22
        \item 28 et 49
        \item 8 et 10
        \item 12 et 30
      \end{enumerate}
    \end{exercise}
    
    \begin{exercise}
      Calcule le PPCM des nombres suivants.
      \begin{enumerate}[cols=3]
        \item 3 et 9
        \item 6 et 12
        \item 25 et 35
        \item 8 et 10
        \item 12 et 30
        \item 14 et 21
      \end{enumerate}
    \end{exercise}
    
    \begin{exercise}
        Écris l'expression littérale générale qui correspond aux énoncés suivants (avec \(n\) pour représenter les naturels).
        \begin{enumerate}[cols=2]
            \item Un nombre pair.
            \item Un multiple de 4 et 8.
            \item La moitié d'un nombre. 
            \item Un nombre impair diminué de 5.
            \item Un nombre impair soustrait de 13.
            \item Un nombre pair multiplié par 4.
        \end{enumerate}
      \end{exercise}


\newpage


\begin{center}
\bfseries
\begin{Huge}
Chapitre 3 - Diviseurs et multiples
\end{Huge}
\smallbreak
\begin{large}
Niveau 1 \faIcon{star} - Fiche 4
\end{large}
\smallbreak
\rule[10pt]{.5\linewidth}{1 pt}
\end{center}

\setcounter{exercise}{0}
    
    \begin{exercise}
      Décompose les nombres suivants en un produit de facteurs premiers.

      \begin{enumerate}[cols=3]
          \item 84
          \item 66
          \item 36
          \item 49
          \item 90
          \item 128
      \end{enumerate}
    \end{exercise}
    
    \begin{exercise}
      Écris l'égalité euclidienne qui découle de cette
    division.
      \begin{enumerate}[cols=3]
        \item $39 : 7$
        \item $42 : 6$
        \item $17 : 4$
        \item $23 : 5$
        \item $11 : 3$
        \item $14 : 6$
      \end{enumerate}
    \end{exercise}
    
    \begin{exercise}
      Calcule le PGCD des nombres suivants.
    \begin{enumerate}[cols=3]
        \item 14 et 21
        \item 36 et 48
        \item 45 et 60
        \item 72 et 108
        \item 81 et 99
        \item 63 et 84
      \end{enumerate}
    \end{exercise}
    
    \begin{exercise}
      Calcule le PPCM des nombres suivants.
      \begin{enumerate}[cols=3]
        \item 36 et 48
        \item 11 et 22
        \item 9 et 14
        \item 10 et 15
        \item 21 et 28
        \item 16 et 24
      \end{enumerate}
    \end{exercise}
    
    \begin{exercise}
        Écris l'expression littérale générale qui correspond aux énoncés suivants (avec \(n\) pour représenter les naturels).
        \begin{enumerate}[cols=2]
            \item Deux nombres consécutifs. 
            \item Deux multiples de 7 consécutifs.
            \item La somme de cinq nombres consécutifs.
            \item Un multiple de 4 augmenté de 7.
            \item Un nombre impair diminué de 2.
        \end{enumerate}
      \end{exercise}

\newpage

\begin{center}
  \bfseries
  \begin{Huge}
    Correctif
  \end{Huge}
  \smallbreak
  \begin{large}
    Niveau 1 \faIcon{star} - Fiche 3
  \end{large}
  \smallbreak
  \rule[10pt]{.5\linewidth}{1 pt}
\end{center}

\setcounter{exercise}{0}
    
    \begin{exercise}
      Décompose les nombres suivants en un produit de facteurs premiers.

          \begin{enumerate}[cols=3]
            \item 54 = $2 \cdot 3 \cdot 3 \cdot 3$
            \item 56 = $2 \cdot 2 \cdot 2 \cdot 7$
            \item 72 = $2 \cdot 2 \cdot 2 \cdot 3 \cdot 3$
            \item 90 = $2 \cdot 3 \cdot 3 \cdot 5$
            \item 48 = $2 \cdot 2 \cdot 2 \cdot 2 \cdot 3$
            \item 70 = $2 \cdot 5 \cdot 7$
          \end{enumerate}
    \end{exercise}
    
    
    \begin{exercise}
      Écris l'égalité euclidienne qui découle de cette
    division.
        \begin{enumerate}[cols=3]
          \item $30 = 8 \cdot 3 + 6$
          \item $28 = 7 \cdot 4 + 0$
          \item $35 = 6 \cdot 5 + 5$
          \item $45 = 9 \cdot 5 + 0$
          \item $24 = 5 \cdot 4 + 4$
          \item $22 = 4 \cdot 5 + 2$
        \end{enumerate}
    \end{exercise}
    
    \begin{exercise}
      Calcule le PGCD des nombres suivants.
        \begin{enumerate}[cols=3]
            \item 3
            \item 5
            \item 11
            \item 7
            \item 2
            \item 6
        \end{enumerate}
    \end{exercise}
    
    \begin{exercise}
      Calcule le PPCM des nombres suivants.
        \begin{enumerate}[cols=3]
            \item 9
            \item 12
            \item 175
            \item 40
            \item 60
            \item 42
        \end{enumerate}
    \end{exercise}
    
    \begin{exercise}
        Écris l'expression littérale générale qui correspond aux énoncés suivants (avec \(n\) pour représenter les naturels).
        \begin{enumerate}[cols=2]
            \item 2n
            \item 8n+8
            \item n : 2 ou \(\dfrac{n}{2}\) 
            \item 2n + 1 - 5 
            \item 13 - 2n + 1 
            \item \(4 \cdot 2n = 8n\) 
        \end{enumerate}
      \end{exercise}

\newpage


\begin{center}
  \bfseries
  \begin{Huge}
    Correctif
  \end{Huge}
  \smallbreak
  \begin{large}
    Niveau 1 \faIcon{star} - Fiche 4
  \end{large}
  \smallbreak
  \rule[10pt]{.5\linewidth}{1 pt}
\end{center}

\setcounter{exercise}{0}
      
      \begin{exercise}
        Décompose les nombres suivants en un produit de facteurs premiers.

          \begin{enumerate}[cols=3]
              \item 84 = $2^2 \cdot 3 \cdot 7$
              \item 66 = $2 \cdot 3 \cdot 11$
              \item 36 = $2^2 \cdot 3^2$
              \item 49 = $7^2$
              \item 90 = $2 \cdot 3^2 \cdot 5$
              \item 128 = $2^7$
            \end{enumerate}
      \end{exercise}
      
      
      \begin{exercise}
        Écris l'égalité euclidienne qui découle de cette
      division.
          \begin{enumerate}[cols=3]
            \item $39 = 7 \cdot 5 + 4$
            \item $42 = 6 \cdot 7 + 0$
            \item $17 = 4 \cdot 4 + 1$
            \item $23 = 5 \cdot 4 + 3$
            \item $11 = 3 \cdot 3 + 2$
            \item $14 = 6 \cdot 2 + 2$
          \end{enumerate}
      \end{exercise}
      
      \begin{exercise}
        Calcule le PGCD des nombres suivants.
          \begin{enumerate}[cols=3]
              \item 7
              \item 12
              \item 15
              \item 36
              \item 9
              \item 21
          \end{enumerate}
      \end{exercise}
      
      \begin{exercise}
        Calcule le PPCM des nombres suivants.
          \begin{enumerate}[cols=3]
              \item 144
              \item 22
              \item 126
              \item 30
              \item 84
              \item 48 
          \end{enumerate}
      \end{exercise}
      
      \begin{exercise}
          Écris l'expression littérale générale qui correspond aux énoncés suivants (avec \(n\) pour représenter les naturels).
          \begin{enumerate}[cols=2]
              \item n et n+1  
              \item $7n$ et $7n+7$
              \item $n + n+1 + n+2 + n+3 + n+4 = 5n + 10$
              \item $4n + 7$
              \item $2n +1 - 2 = 2n -1$
          \end{enumerate}
        \end{exercise}


%Fin niveau 1%%%%%%%%%%%%%%%%%%%%%%%%%%%%%%%%%%%%%%%%%%

\newpage

\begin{center}
  \bfseries
  \begin{Huge}
    Chapitre 3 - Diviseurs et multiples
  \end{Huge}
  \smallbreak
  \begin{large}
    Niveau 2 \faIcon{star}\faIcon{star} - Fiche 1
  \end{large}
  \smallbreak
  \rule[10pt]{.5\linewidth}{1 pt}
\end{center}

\setcounter{exercise}{0}
      
      \begin{exercise}
        Décompose les nombres suivants en un produit de facteurs premiers.
          \begin{enumerate}[cols=4]
            \item $180$
            \item $10 920$
            \item $54$
            \item $2340$
          \end{enumerate}
      \end{exercise}
      
      \begin{exercise}
        Effectue la division euclidienne, puis écris l'égalité qui en découle. 
          \begin{enumerate}[cols=2]
            \item $186 : 5$
            \item $8652 : 11$
            \item $\dfrac{657}{3}$
            \item $\dfrac{8756}{20}$
          \end{enumerate}
      \end{exercise}
      
      \begin{exercise}
        Calcule le PGCD des nombres suivants. 
          \begin{enumerate}[cols=2]
            \item $2400, 3600$ et $48000$
            \item $84$ et $12$
            \item $75, 125$ et $300$
            \item $14$ et $15$
          \end{enumerate}
      \end{exercise}
      
      \begin{exercise}
        Calcule le PPCM des nombres suivants.
          \begin{enumerate}[cols=2]
            \item $40$ et $16$
            \item $75$ et $90$
            \item $90$ et $48$
            \item $7$ et $15$
          \end{enumerate}
      \end{exercise}
      
      \begin{exercise}
          Prouve, à l'aide des expressions littérales (\(n\) représentant n'importe quel naturel) \dots 
          \begin{enumerate}
            \item qu'un multiple de 15 est multiple de 5
            \item qu'un multiple de 40 est multiple de 5 
            \item que 0 est un nombre pair
          \end{enumerate}
        \end{exercise}

\newpage

  \begin{center}
    \bfseries
    \begin{Huge}
      Chapitre 3 - Diviseurs et multiples
    \end{Huge}
    \smallbreak
    \begin{large}
      Niveau 2 \faIcon{star}\faIcon{star} - Fiche 2
    \end{large}
    \smallbreak
    \rule[10pt]{.5\linewidth}{1 pt}
  \end{center}
  
  \setcounter{exercise}{0}
        
        \begin{exercise}
          Décompose les nombres suivants en un produit de facteurs premiers.
            \begin{enumerate}[cols=2]
              \item $8820$
              \item $48$
              \item $1650$
              \item $2048$
            \end{enumerate}
        \end{exercise}
        
        \begin{exercise}
          Effectue la division euclidienne, puis écris l'égalité qui en découle. 
            \begin{enumerate}[cols=2]
              \item $\dfrac{86}{6}$
              \item $946 : 12$
              \item $\dfrac{308}{6}$
              \item $1685 : 5$
            \end{enumerate}
        \end{exercise}
 
        \begin{exercise}
          Calcule le PGCD des nombres suivants. 
            \begin{enumerate}[cols=2]
              \item $9,30$ et $15$
              \item $100$ et $50$
              \item $625$ et $650$
              \item $360, 540$ et $720$
            \end{enumerate}
        \end{exercise}
 
        \begin{exercise}
          Calcule le PPCM des nombres suivants.
            \begin{enumerate}[cols=2]
              \item $80$ et $84$
              \item $44$ et $55$
              \item $150$ et $300$
              \item $2$ et $3$
            \end{enumerate}
        \end{exercise}

        \begin{exercise}
          Prouve, à l'aide des expressions littérales (\(n\) représentant n'importe quel naturel) \dots 
          \begin{enumerate}
            \item qu'un multiple de 21 est multiple de 3
            \item qu'un multiple de 40 est multiple de 2 
            \item que 18n + 9 n'est pas un multiple de 18
          \end{enumerate}
        \end{exercise}

\newpage

\begin{center}
  \bfseries
  \begin{Huge}
    Correctif
  \end{Huge}
  \smallbreak
  \begin{large}
    Niveau 2 \faIcon{star}\faIcon{star} - Fiche 1
  \end{large}
  \smallbreak
  \rule[10pt]{.5\linewidth}{1 pt}
\end{center}

\setcounter{exercise}{0}
      
      \begin{exercise}
        Décompose les nombres suivants en un produit de facteurs premiers.
          \begin{enumerate}[cols=2]
            \item $180 = 2^2 \cdot 3^2 \cdot 5$
            \item $10 920 = 2^3 \cdot 3^2 \cdot 5 \cdot 7$
            \item $54 = 2 \cdot 3^3$
            \item $2340 = 2^2 \cdot 3 \cdot 5 \cdot 13$
          \end{enumerate}
      \end{exercise}
      
      \begin{exercise} 
          \begin{enumerate}[cols=2]
            \item $186=5\cdot37+1$
            \item $8652=11\cdot786+6$
            \item $657=3\cdot219 (+0)$
            \item $8756=20\cdot437+16$
          \end{enumerate}
      \end{exercise}
      
      \begin{exercise}
        Calcule le PGCD des nombres suivants. 
          \begin{enumerate}[cols=2]
            \item  $1200$
            \item  $12$
            \item $25$
            \item $1$
          \end{enumerate}
      \end{exercise}
      
      \begin{exercise}
        Calcule le PPCM des nombres suivants.
          \begin{enumerate}[cols=2]
            \item $80$
            \item $450$
            \item $720$
            \item $105$
          \end{enumerate}
      \end{exercise}
      
      \begin{exercise}
          Prouve, à l'aide des expressions littérales (\(n\) représentant n'importe quel naturel) \dots 
          \begin{enumerate}
            \item \(15n = 5 \cdot 3n\)
            \item \(40n = 5 \cdot 8n\)
            \item \(2 \cdot 0 = 0\) 
          \end{enumerate}
        \end{exercise}


\newpage
  \begin{center}
  \bfseries
  \begin{Huge}
  Correctif
  \end{Huge}
  \smallbreak
  \begin{large}
  Niveau 2 \faIcon{star}\faIcon{star} - Fiche 2
  \end{large}
  \smallbreak
  \rule[10pt]{.5\linewidth}{1 pt}
  \end{center}

\setcounter{exercise}{0}
  
  \begin{exercise}
    Décompose les nombres suivants en un produit de facteurs premiers.
      \begin{enumerate}[cols=2]
        \item $8820 = 2^2 \cdot 3^2 \cdot 5 \cdot 7^2$
        \item $48 = 2^4 \cdot 3$
        \item $1650 = 2 \cdot 3 \cdot 5^2 \cdot 11$
        \item $2048 = 2^{11}$
      \end{enumerate}
  \end{exercise}
  
  \begin{exercise}
      \begin{enumerate}[cols=2]
        \item $86=6\cdot14+2$
        \item $946=12\cdot78+10$
        \item $308=6\cdot51+2$
        \item $1685=5\cdot337+0$
      \end{enumerate}
  \end{exercise}
  
  \begin{exercise}
    Calcule le PGCD des nombres suivants. 
      \begin{enumerate}[cols=2]
        \item $3$
        \item $50$
        \item $25$
        \item $180$
      \end{enumerate}
  \end{exercise}
  
  \begin{exercise}
    Calcule le PPCM des nombres suivants.
      \begin{enumerate}[cols=2]
        \item $1680$
        \item $220$
        \item $300$
        \item $6$
      \end{enumerate}
  \end{exercise}
  
  \begin{exercise}
    Prouve, à l'aide des expressions littérales (\(n\) représentant n'importe quel naturel) \dots 
    \begin{enumerate}
      \item \(21n = 3 \cdot 7n\)
      \item \(40n = 2 \cdot 20n\)
      \item si n=1, alors \(18\cdot 1 + 9 = 27\) et 27 n'est pas un multiple de 18
    \end{enumerate}
  \end{exercise}
      
\newpage


\begin{center}
  \bfseries
  \begin{Huge}
    Chapitre 3 - Diviseurs et multiples
  \end{Huge}
  \smallbreak
  \begin{large}
    Niveau 2 \faIcon{star}\faIcon{star} - Fiche 3
  \end{large}
  \smallbreak
  \rule[10pt]{.5\linewidth}{1 pt}
\end{center}

\setcounter{exercise}{0}
      
      \begin{exercise}
        Décompose les nombres suivants en un produit de facteurs premiers.
          \begin{enumerate}[cols=2]
            \item $91$
            \item $144$
            \item $1875$
            \item $1650$
          \end{enumerate}
      \end{exercise}
      
      \begin{exercise}
        Effectue la division euclidienne, puis écris l'égalité qui en découle. 
          \begin{enumerate}[cols=2]
            \item $\dfrac{242424}{12}$
            \item $9063 : 7 $
            \item $658 : 12$
            \item $\dfrac{1385}{15}$
          \end{enumerate}
      \end{exercise}
      
      \begin{exercise}
        Calcule le PGCD des nombres suivants. 
          \begin{enumerate}[cols=2]
            \item $81$ et $540$
            \item $18, 36$ et $66$
            \item $144$ et $288$
            \item $17$ et $23$
          \end{enumerate}
      \end{exercise}
      
      \begin{exercise}
        Calcule le PPCM des nombres suivants.
          \begin{enumerate}[cols=2]
            \item $12$ et $15$
            \item $50$ et $75$
            \item $180$ et $270$
            \item $15$ et $75$
          \end{enumerate}
      \end{exercise}
      
      \begin{exercise}
        Prouve, à l'aide des expressions littérales (\(n\) représentant n'importe quel naturel) \dots 
        \begin{enumerate}
          \item qu'un multiple de 15 augmenté de 6 est un multiple de 3
          \item qu'un multiple de 70 est multiple de 5 
          \item que \(5n+10\) n'est pas un nombre pair
        \end{enumerate}
      \end{exercise}

\newpage


\begin{center}
  \bfseries
  \begin{Huge}
    Chapitre 3 - Diviseurs et multiples
  \end{Huge}
  \smallbreak
  \begin{large}
    Niveau 2 \faIcon{star}\faIcon{star} - Fiche 4
  \end{large}
  \smallbreak
  \rule[10pt]{.5\linewidth}{1 pt}
\end{center}

\setcounter{exercise}{0}
      
      \begin{exercise}
        Décompose les nombres suivants en un produit de facteurs premiers.
          \begin{enumerate}[cols=2]
            \item $120$
            \item $13 860$
            \item $630$
            \item $64$
          \end{enumerate}
      \end{exercise}

      \begin{exercise}
        Effectue la division euclidienne, puis écris l'égalité qui en découle. 
          \begin{enumerate}[cols=2]
            \item $298 : 18$
            \item $\dfrac{527}{24} $
            \item $\dfrac{724}{8}$
            \item $938 : 12$
          \end{enumerate}
      \end{exercise}

      \begin{exercise}
        Calcule le PGCD des nombres suivants. 
          \begin{enumerate}[cols=2]
            \item $180$ et $168$
            \item $900$ et $360$
            \item $81$ et $428$
            \item $1125, 1575$ et $2025$
          \end{enumerate}
      \end{exercise}

      \begin{exercise}
        Calcule le PPCM des nombres suivants.
          \begin{enumerate}[cols=2]
            \item $16$ et $12$
            \item $14$ et $10$
            \item $60$ et $35$
            \item $300$ et $27$
          \end{enumerate}
      \end{exercise}

      \begin{exercise}
        Prouve, à l'aide des expressions littérales (\(n\) représentant n'importe quel naturel) \dots 
        \begin{enumerate}
          \item que \(3n + 27\) n'est pas un multiple de 5
          \item qu'un multiple de 22 est mulitple 11 
          \item que \(36n + 20\) est un multiple de 4
        \end{enumerate}
      \end{exercise}

\newpage
\begin{center}
  \bfseries
  \begin{Huge}
    Correctif
  \end{Huge}
  \smallbreak
  \begin{large}
    Niveau 2 \faIcon{star}\faIcon{star} - Fiche 3
  \end{large}
  \smallbreak
  \rule[10pt]{.5\linewidth}{1 pt}
\end{center}

\setcounter{exercise}{0}
      
      \begin{exercise}
        Décompose les nombres suivants en un produit de facteurs premiers.
          \begin{enumerate}[cols=2]
            \item $91 = 7 \cdot 13$
            \item $144 = 2^4 \cdot 3^2$
            \item $1875 = 3^1 \cdot 5^3 \cdot 5^3$
            \item $1650 = 2 \cdot 3 \cdot 5^2 \cdot 11$
          \end{enumerate}
      \end{exercise}
      
      \begin{exercise}
        Effectue la division euclidienne, puis écris l'égalité qui en découle. 
          \begin{enumerate}[cols=2]
            \item $242424=12\cdot20202+0$
            \item $9063=7\cdot1294+5$
            \item $658=12\cdot54+10$
            \item $1385=15\cdot92+5$
          \end{enumerate}
      \end{exercise}
      
      \begin{exercise}
        Calcule le PGCD des nombres suivants.
          \begin{enumerate}[cols=2]
            \item $27$
            \item $6$
            \item $144$
            \item $1$
          \end{enumerate}
      \end{exercise}
      
      \begin{exercise}
        Calcule le PPCM des nombres suivants.
          \begin{enumerate}[cols=2]
            \item $60$
            \item $150$
            \item $54$
            \item $75$
          \end{enumerate}
      \end{exercise}
      
      \begin{exercise}
        Prouve, à l'aide des expressions littérales (\(n\) représentant n'importe quel naturel) \dots 
        \begin{enumerate}
          \item \(15n + 6 = 3 \cdot (5n+2)\)
          \item \(70n = 5 \cdot 14n\)
          \item si n=3, alors \(5\cdot 3 + 10 = 25\) et 25 n'est pas un nombre pair
        \end{enumerate}
      \end{exercise}
      
\newpage

\begin{center}
  \bfseries
  \begin{Huge}
  Correctif
  \end{Huge}
  \smallbreak
  \begin{large}
  Niveau 2 \faIcon{star}\faIcon{star} - Fiche 4
  \end{large}
  \smallbreak
  \rule[10pt]{.5\linewidth}{1 pt}
  \end{center}

\setcounter{exercise}{0}
  
  \begin{exercise}
    Décompose les nombres suivants en un produit de facteurs premiers.
      \begin{enumerate}[cols=2]
        \item $120 = 2^3 \cdot  3 \cdot 5$
        \item $13 860 = 2^2 \cdot 3^3 \cdot 5 \cdot 7^2$
        \item $630 = 2 \cdot 3^2 \cdot 5 \cdot 7$
        \item $64 = 2^6$
      \end{enumerate}
  \end{exercise}
  
  \begin{exercise}
    Effectue la division euclidienne, puis écris l'égalité qui en découle. 
      \begin{enumerate}[cols=2]
        \item $298 = 18 \cdot 16 + 10$
        \item $2527 = 24 \cdot 21 + 23$
        \item $724 = 8 \cdot 90 + 4$
        \item $ 938 = 12 \cdot 78 + 2$
      \end{enumerate}
  \end{exercise}
  
  \begin{exercise}
    Calcule le PGCD des nombres suivants. 
      \begin{enumerate}[cols=2]
        \item $12$
        \item $180$
        \item $1$
        \item $225$
      \end{enumerate}
  \end{exercise}
  
  \begin{exercise}
    Calcule le PPCM des nombres suivants.
      \begin{enumerate}[cols=2]
        \item $48$
        \item $70$
        \item $420$
        \item $2700$
      \end{enumerate}
  \end{exercise}
  
  \begin{exercise}
    Prouve, à l'aide des expressions littérales (\(n\) représentant n'importe quel naturel) \dots 
    \begin{enumerate}
      \item si n = 2, alors \(3 \cdot 2 + 27 = 33\) et 33 n'est pas un multiple 5 
      \item \(22n = 11 \cdot 2n\) 
      \item \(36n+20 = 4 (9n+5)\)
    \end{enumerate}
  \end{exercise}
      
  
  % Fin niveau 2 %%%%%%%%%%%%%%%%%%%%%%%%%%%%%%%%%%%%%%%

\newpage

\begin{center}
  \bfseries
  \begin{Huge}
    Chapitre 3 - Diviseurs et multiples
  \end{Huge}
  \smallbreak
  \begin{large}
    Niveau 3 \faIcon{star} \faIcon{star} \faIcon{star} - Fiche 1
  \end{large}
  \smallbreak
  \rule[10pt]{.5\linewidth}{1 pt}
\end{center}

\setcounter{exercise}{0}
      
\begin{exercise}
  Retrouve la propriété qui permet d'affirmer, sans calculer, que... 
  \begin{enumerate}[cols=2]
    \item PGCD(700, 13)= 1
    \item PPCM(350, 175)= 350
    \item PPCM(52, 9)= 468
    \item PGCD (15, 45)= 15
  \end{enumerate}
\end{exercise}

\begin{exercise}
  Effectue la division euclidienne, puis écris l'égalité qui en découle. 
    \begin{enumerate}[cols=2]
      \item \(1385 : 15\)
      \item \(658 : 12\)
      \item \(\dfrac{9063}{7}\)
      \item \(3801 : 3\)
    \end{enumerate}
\end{exercise}

\begin{exercise}
  Calcule le PGCD et le PPCM des nombres suivants. 
    \begin{enumerate}[cols=2]
      \item 12, 14 et 6
      \item 14 et 15
      \item 120, 36 et 200
      \item 45, 600 et 90
    \end{enumerate}
\end{exercise}

\newpage

\begin{center}
  \bfseries
  \begin{Huge}
    Chapitre 3 - Diviseurs et multiples
  \end{Huge}
  \smallbreak
  \begin{large}
    Niveau 3 \faIcon{star} \faIcon{star} \faIcon{star} - Fiche 2
  \end{large}
  \smallbreak
  \rule[10pt]{.5\linewidth}{1 pt}
\end{center}

\setcounter{exercise}{0}
      
\begin{exercise}
  Retrouve la propriété qui permet d'affirmer, sans calculer, que... 
  \begin{enumerate}[cols=2]
    \item PGCD(24, 55)= 1
    \item PPCM(14, 70)= 70
    \item PPCM(2, 25)= 50
    \item PGCD (500, 1000)= 500
  \end{enumerate}
\end{exercise}

\begin{exercise}
  Effectue la division euclidienne, puis écris l'égalité qui en découle. 
    \begin{enumerate}[cols=2]
      \item \(186 : 5\)
      \item \(8652 : 11\)
      \item \(\dfrac{8756}{11}\)
      \item \(3579 : 13\)
    \end{enumerate}
\end{exercise}

\begin{exercise}
  Calcule le PGCD et le PPCM des nombres suivants. 
    \begin{enumerate}[cols=2]
      \item 80, 90 et 100
      \item 25, 45 et 100
      \item 1800, 900 et 90
      \item 27, 15 et 45
    \end{enumerate}
\end{exercise}

\newpage

\begin{center}
  \bfseries
  \begin{Huge}
    Correctif
  \end{Huge}
  \smallbreak
  \begin{large}
    Niveau 3 \faIcon{star} \faIcon{star} \faIcon{star} - Fiche 1
  \end{large}
  \smallbreak
  \rule[10pt]{.5\linewidth}{1 pt}
\end{center}

\setcounter{exercise}{0}
      
\begin{exercise}
  Retrouve la propriété qui permet d'affirmer, sans calculer, que... 
  \begin{enumerate}
    \item Si deux nombres sont premiers entre eux, alors leur PGCD vaut 1
    \item Si deux nombres sont multiples l'un de l'autre, leur PPCM vaut le plus grand des deux nombres
    \item Si deux nombres sont premiers entre eux, alors leur PPCM vaut leur produit
    \item Si deux nombres sont multiples l'un de l'autre, leur PGCD vaut le plus petit des deux nombres
  \end{enumerate}
\end{exercise}

\begin{exercise}
  Effectue la division euclidienne, puis écris l'égalité qui en découle.
    \begin{enumerate}[cols=2]
      \item \(1385 = 15 \cdot 92 + 5\)
      \item \(658 = 12 \cdot 54 + 10\)
      \item \(9063=7 \cdot 1294 + 5\)
      \item \(3801 = 3 \cdot 1267 (+0)\)
    \end{enumerate}
\end{exercise}

\begin{exercise}
  Calcule le PGCD et le PPCM des nombres suivants.
    \begin{enumerate}[cols=2]
      \item  PGCD(12, 14, 6) = 2 \\ PPCM(12, 14, 6) = 84 
      \item PGCD (14, 15) = 1 \\ PPCM(14, 15)= 70
      \item PGCD(120, 36, 200) = 4 \\ PPCM(120, 36, 200)= 1800
      \item PGCD(45, 600, 90) = 15 \\ PPCM(45, 600, 90) = 1800
    \end{enumerate}
\end{exercise}


\newpage

\begin{center}
  \bfseries
  \begin{Huge}
    Correctif
  \end{Huge}
  \smallbreak
  \begin{large}
    Niveau 3 \faIcon{star} \faIcon{star} \faIcon{star} - Fiche 2
  \end{large}
  \smallbreak
  \rule[10pt]{.5\linewidth}{1 pt}
\end{center}

\setcounter{exercise}{0}
      
\begin{exercise}
  Retrouve la propriété qui permet d'affirmer, sans calculer, que... 
  \begin{enumerate}
    \item Si deux nombres sont premiers entre eux, alors leur PGCD vaut 1
    \item Si deux nombres sont multiples l'un de l'autre, leur PPCM vaut le plus grand des deux nombres
    \item Si deux nombres sont premiers entre eux, alors leur PPCM vaut leur produit
    \item Si deux nombres sont multiples l'un de l'autre, leur PGCD vaut le plus petit des deux nombres
  \end{enumerate}
\end{exercise}

\begin{exercise}
  Effectue la division euclidienne, puis écris l'égalité qui en découle. 
    \begin{enumerate}[cols=2]
      \item \(186=5\cdot37 +1\)
      \item \(8652 = 11 \cdot 786 +6\)
      \item \(8756=11 \cdot 796 (+0)\)
      \item \(3579 = 13 \cdot 275 + 4\)
    \end{enumerate}
\end{exercise}

\begin{exercise}
  Calcule le PGCD et le PPCM des nombres suivants. 
    \begin{enumerate}[cols=2]
      \item  PGCD(80, 90, 100) = 10 \\ PPCM(80, 90, 100) = 3600 
      \item PGCD (25, 45 et 100) = 5 \\ PPCM(25, 45 et 100)= 900
      \item PGCD(1800, 900 et 90) = 90 \\ PPCM(1800, 900 et 90)= 1800
      \item PGCD(27, 15 et 45) = 3 \\ PPCM(27, 15 et 45) = 135
    \end{enumerate}
\end{exercise}

\newpage

\begin{center}
  \bfseries
  \begin{Huge}
    Chapitre 3 - Diviseurs et multiples
  \end{Huge}
  \smallbreak
  \begin{large}
    Niveau 3 \faIcon{star} \faIcon{star} \faIcon{star} - Fiche 3
  \end{large}
  \smallbreak
  \rule[10pt]{.5\linewidth}{1 pt}
\end{center}

\setcounter{exercise}{0}
      
\begin{exercise}
  Retrouve la propriété qui permet d'affirmer, sans calculer, que... 
  \begin{enumerate}[cols=2]
    \item PGCD(35, 140)= 35
    \item PGCD(1, 15)= 1
    \item PGCD(60, 19)= 1
    \item PGCD (8, 40)= 8
  \end{enumerate}
\end{exercise}

\begin{exercise}
  Effectue la division euclidienne, puis écris l'égalité qui en découle. 
    \begin{enumerate}[cols=2]
      \item \(3298 : 12\)
      \item \(4002 : 5\)
      \item \(\dfrac{2139}{20}\)
      \item \(1408:2\)
    \end{enumerate}
\end{exercise}

\begin{exercise}
  Calcule le PGCD et le PPCM des nombres suivants. 
    \begin{enumerate}[cols=2]
      \item 48 et 200
      \item 108 et 500
      \item 66 et 140
      \item 270 et 84
    \end{enumerate}
\end{exercise}

\newpage

\begin{center}
  \bfseries
  \begin{Huge}
    Chapitre 3 - Diviseurs et multiples
  \end{Huge}
  \smallbreak
  \begin{large}
    Niveau 3 \faIcon{star} \faIcon{star} \faIcon{star} - Fiche 4
  \end{large}
  \smallbreak
  \rule[10pt]{.5\linewidth}{1 pt}
\end{center}

\setcounter{exercise}{0}
      
\begin{exercise}
  Retrouve la propriété qui permet d'affirmer, sans calculer, que... 
  \begin{enumerate}[cols=2]
    \item PGCD(45, 8)= 1
    \item PPCM(200, 600)= 600
    \item PPCM(25, 24)= 600
    \item PGCD (46, 23)= 23
  \end{enumerate}
\end{exercise}

\begin{exercise}
  Effectue la division euclidienne, puis écris l'égalité qui en découle. 
    \begin{enumerate}[cols=2]
      \item \(132 : 11\)
      \item \(131 : 12\)
      \item \(\dfrac{2703}{15}\)
      \item \(3801 : 3\)
    \end{enumerate}
\end{exercise}

\begin{exercise}
  Calcule le PGCD et le PPCM des nombres suivants. 
    \begin{enumerate}[cols=2]
      \item 21, 54 et 12
      \item 5 et 12
      \item 540, 360 et 180
      \item 22, 35 et 55
    \end{enumerate}
\end{exercise}

\newpage

\begin{center}
  \bfseries
  \begin{Huge}
    Correctif
  \end{Huge}
  \smallbreak
  \begin{large}
    Niveau 3 \faIcon{star} \faIcon{star} \faIcon{star} - Fiche 3
  \end{large}
  \smallbreak
  \rule[10pt]{.5\linewidth}{1 pt}
\end{center}

\setcounter{exercise}{0}
      
\begin{exercise}
  Retrouve la propriété qui permet d'affirmer, sans calculer, que... 
  \begin{enumerate}
    \item Si deux nombres sont multiples l'un de l'autre, leur PGCD vaut le plus petit des deux nombres
    \item Si deux nombres sont multiples l'un de l'autre, leur PGCD vaut le plus petit des deux nombres
    \item Si deux nombres sont premiers entre eux, alors leur PGCD vaut 1
    \item Si deux nombres sont multiples l'un de l'autre, leur PGCD vaut le plus petit des deux nombres
  \end{enumerate}
\end{exercise}

\begin{exercise}
Effectue la division euclidienne, puis écris l'égalité qui en découle. 
\begin{enumerate}[cols=2]
\item \(3298 = 12 \cdot 274 + 10\)
\item \(4002 = 5 \cdot 800 + 2\)
\item \(2139=20 \cdot 106 + 19\)
\item \(1408 = 2 \cdot 704 (+0)\)
\end{enumerate}
\end{exercise}

\begin{exercise}
Calcule le PGCD et le PPCM des nombres suivants. 
\begin{enumerate}[cols=2]
\item  PGCD(48, 200) = 8 \\ PPCM(48, 200) = 1200 
\item PGCD (108, 500) = 4 \\ PPCM(108, 500)= 13500
\item PGCD(66, 140) = 2 \\ PPCM(66, 140)= 4620
\item PGCD(270, 84) = 6 \\ PPCM(270, 84) = 3780
\end{enumerate}
\end{exercise}

\newpage

\begin{center}
  \bfseries
  \begin{Huge}
    Correctif
  \end{Huge}
  \smallbreak
  \begin{large}
    Niveau 3 \faIcon{star} \faIcon{star} \faIcon{star} - Fiche 4
  \end{large}
  \smallbreak
  \rule[10pt]{.5\linewidth}{1 pt}
\end{center}

\setcounter{exercise}{0}
      
\begin{exercise}
  Retrouve la propriété qui permet d'affirmer, sans calculer, que...
  \begin{enumerate}
    \item Si deux nombres sont premiers entre eux, alors leur PGCD vaut 1
    \item Si deux nombres sont multiples l'un de l'autre, leur PPCM vaut le plus grand des deux nombres
    \item Si deux nombres sont premiers entre eux, alors leur PPCM vaut leur produit
    \item Si deux nombres sont multiples l'un de l'autre, leur PGCD vaut le plus petit des deux nombres
  \end{enumerate}
\end{exercise}

\begin{exercise}
  Effectue la division euclidienne, puis écris l'égalité qui en découle.
    \begin{enumerate}[cols=2]
      \item \(132 = 11 \cdot 12 (+ 0)\)
      \item \(131 = 12 \cdot 10 + 11\)
      \item \(2703=15 \cdot 180 + 3\)
      \item \(3801 = 3 \cdot 1267 (+0)\)
    \end{enumerate}
\end{exercise}

\begin{exercise}
  Calcule le PGCD et le PPCM des nombres suivants. 
    \begin{enumerate}[cols=2]
      \item  PGCD(21, 54, 12) = 3 \\ PPCM(21, 54, 12) = 756 
      \item PGCD (5, 12) = 1 \\ PPCM(5, 12)= 60
      \item PGCD(540, 360, 180) = 180 \\ PPCM(540, 360, 180)= 1080
      \item PGCD(22, 35, 55) = 1\\ PPCM(22, 35, 55) = 770
    \end{enumerate}
\end{exercise}
  

\end{document}